\section{Coupled Simulation Formulation}
\label{sec:formulation}

\subsection{State and Frames}

\mjorbit{} represents the chief orbit in Earth-centered inertial coordinates
as
\begin{equation}
  x_o = (\bm R,\bm V,t), \qquad \bm R,\bm V \in \R^3,
\end{equation}
with distances in kilometers and time in seconds. The MuJoCo model state
$(q,v)$ is expressed in a local LVLH frame attached to the chief. The LVLH
basis is
\begin{equation}
  \hat{\bm x}_L = \frac{\bm R}{\|\bm R\|}, \qquad
  \hat{\bm z}_L = \frac{\bm R \times \bm V}{\|\bm R \times \bm V\|}, \qquad
  \hat{\bm y}_L = \hat{\bm z}_L \times \hat{\bm x}_L ,
\end{equation}
and the rows of $C_{LI}$ map ECI vectors into LVLH coordinates. The angular
velocity of the LVLH frame is
\begin{equation}
  \bm\omega_L = C_{LI}\frac{\bm R \times \bm V}{\|\bm R\|^2}.
\end{equation}
When J2 is enabled, $\dot{\bm\omega}_L$ is computed from the current
perturbed orbital acceleration rather than assuming two-body motion.

The MuJoCo world frame is identified with the LVLH frame. For each MuJoCo
body $i$, the body center-of-mass position $\bm r_i$ and velocity
$\bm v_i$ are read from MuJoCo in SI units and converted to kilometers where
the orbital layer requires them.

\subsection{Per-Body Orbital Forcing}

The translational orbital coupling applies the exact rotating-frame apparent
acceleration to each body center of mass:
\begin{align}
  \bm a_i^L
    &= C_{LI}\left[\bm g(\bm R_i)-\bm g(\bm R)\right]
       -2\bm\omega_L \times \bm v_i^L \nonumber\\
    &\quad -\dot{\bm\omega}_L \times \bm r_i^L
       -\bm\omega_L \times (\bm\omega_L \times \bm r_i^L),
  \label{eq:apparent-acceleration}
\end{align}
where $\bm R_i = \bm R + C_{IL}\bm r_i^L$ is the absolute ECI position of
body $i$ and $\bm g(\cdot)$ includes point-mass gravity and, optionally, J2.
Equation~\eqref{eq:apparent-acceleration} is not a Clohessy-Wiltshire
linearization, though it reduces to the CW limit for small offsets in
circular two-body orbits. The force applied to body $i$ is
\begin{equation}
  \bm F_i^L = m_i \bm a_i^L .
\end{equation}

\subsection{Environmental Wrenches}

Additional external wrenches are accumulated in a per-body buffer before
being written to MuJoCo's \texttt{xfrc\_applied}. The current implementation
includes flat-plate drag and solar radiation pressure, residual magnetic
dipoles, and rotational gravity-gradient torque. For a configured surface
with body-frame normal $\hat{\bm n}$, center of pressure $\bm p$, and area
$A$, the surface module rotates the normal and center of pressure into LVLH,
computes atmosphere-relative velocity at the point, and applies force and
moment
\begin{equation}
  \bm\tau_i^L = \bm p_i^L \times \bm F_i^L.
\end{equation}
The gravity-gradient torque on a rigid body is modeled as
\begin{equation}
  \bm\tau_{gg}
    = \frac{3\mu}{\|\bm R_i\|^3}
      \hat{\bm r}_i \times \left(J_i^L \hat{\bm r}_i\right),
  \label{eq:gravity-gradient}
\end{equation}
where $J_i^L$ is the body's inertia tensor expressed in the MuJoCo world/LVLH
frame.

\subsection{Actuators and Sensors}

Actuator metadata is compiled against MuJoCo body names, while runtime
commands live on \texttt{MjoData.actuators}. Reaction wheels integrate wheel
speed from commanded torque and apply the equal-and-opposite body torque.
They also apply gyrostat coupling from stored wheel momentum. Magnetorquers
apply $\bm\tau = \bm m \times \bm B$ using the environment magnetic field.
Thrusters apply a clipped body-frame force at a body-fixed point and therefore
produce both force and torque. Single-gimbal control moment gyros are modeled
with a body-fixed gimbal axis $\hat{\bm g}$, spin axis
$\hat{\bm s}(0)$, rotor momentum $h$, and gimbal angle $\theta$:
\begin{equation}
  \hat{\bm s}(\theta)
    = \cos\theta\,\hat{\bm s}(0)
      + \sin\theta\,\left(\hat{\bm g}\times\hat{\bm s}(0)\right),
\end{equation}
with commanded output torque
\begin{equation}
  \bm\tau_{cmg}
    = -h\dot\theta\,
      \left(\hat{\bm g}\times\hat{\bm s}(\theta)\right).
\end{equation}

Sensors are resolved from the MuJoCo sensor catalog and exposed through a
runtime namespace. The package currently supports canonical MuJoCo sensor
data access, stochastic measurement helpers, and custom orbit-aware sensors
for quantities such as sun, horizon, star, and magnetic-field directions.

\subsection{Stepping Algorithm}

Each call to \texttt{mjo\_step} clears the external-wrench buffer, assembles
environmental and actuator wrenches, copies them into MuJoCo, computes the
net translational external force for chief-orbit feedback, advances the
chief orbit with RK4, refreshes frame and environment caches, and finally
steps the MuJoCo state. Internal MuJoCo forces such as joints, constraints,
and contacts are intentionally not fed back into the chief orbit; only net
external forces accumulated in the coupling layer affect the propagated
chief state.
